\chapter{Estado del arte}
\label{ch:estado-del-arte}

La conducción autónoma se define como la capacidad de un vehículo para
navegar de manera segura y eficiente prescindiendo de la intervención
directa de un operador humano~\cite{ita_deteccion}. Esta tecnología se
ha consolidado como una de las más disruptivas para el futuro de la
movilidad, centrando su potencial en la mejora de la seguridad vial, la
optimización del transporte y la accesibilidad.

No obstante, el concepto de automatización vehicular no es reciente.
Sus orígenes se remontan a siglos atrás, con ideas y prototipos que han
evolucionado hasta dar lugar a los avanzados sistemas actuales. Para
situar el presente proyecto en su contexto tecnológico es necesario
analizar la evolución histórica de la conducción autónoma, los
distintos niveles de automatización definidos por organismos
especializados y las principales tecnologías que hacen posible su
funcionamiento.

\section{Historia de los vehículos autónomos}
\label{sec:historia-va}

La historia de los vehículos autónomos comenzó a finales del siglo~XV
cuando Leonardo Da Vinci diseñó un carro autopropulsado al que se le
podía programar un recorrido predeterminado~\cite{motorpedia_historia}.
Este diseño introdujo, de forma temprana, la idea de un vehículo capaz
de desplazarse sin intervención humana directa.

Posteriormente, la invención del giroscopio por Léon Foucault en 1852
permitió mantener una dirección estable incluso cuando el entorno
cambiaba. A finales del mismo siglo, Ludwig Obry aplicó este principio
en sistemas de navegación autónoma para vehículos
autopropulsados~\cite{prensa_conduccion}.

La automatización vehicular moderna se empezó a conceptualizar en la
Exposición Universal de Nueva York de 1939. En el pabellón
\textit{Futurama}, el diseñador industrial Norman Bel Geddes planteó
que la inteligencia de la navegación debía residir en la
infraestructura viaria, mediante el uso de raíles o circuitos
magnéticos integrados, y no en el propio
vehículo~\cite{geddes1940magic}. Bajo este modelo, el conductor
mantendría el control manual en entornos urbanos y delegaría la
conducción al sistema automático únicamente al incorporarse a
autopistas específicamente equipadas para ello.

Este paradigma cambió con los trabajos del ingeniero alemán Ernst
Dickmanns, quien defendió que el vehículo debía ser capaz de \emph{ver}
y procesar su entorno de forma local, en línea con el enfoque actual de
la robótica móvil~\cite{quintessence_dickmanns}. Logró convertir una
furgoneta Mercedes-Benz en un vehículo autónomo guiado por visión y un
ordenador integrado.

Además, en 1995, gracias a la reducción de tamaño de los ordenadores y
al aumento de sus prestaciones, consiguió que un Mercedes-Benz se
desplazara entre Múnich y Copenhague, casi 1.000~km, con un porcentaje
de conducción autónoma del 95\,\%~\cite{xataka_tecnologia}. A partir de
estos avances, se sentaron las bases de la conducción autónoma
contemporánea.

El punto de inflexión se produjo a principios de este siglo, cuando la
Agencia de Proyectos de Investigación Avanzados del Departamento de
Defensa de Estados Unidos (\textit{Defense Advanced Research Projects
Agency}, DARPA), organizó una serie de competiciones denominadas
\textit{Grand Challenges}~\cite{iptc_vehiculo}. En el primero de ellos
(2004), ningún vehículo logró completar el recorrido de 240~km por el
desierto de Mojave; el mejor participante apenas recorrió 11~km antes
de quedar bloqueado~\cite{elmundo_carrera}. El problema principal fue
la incapacidad de los sistemas para distinguir entre sombras, texturas
de suelo y obstáculos reales, lo que impulsó el desarrollo de nuevas
técnicas de percepción y navegación autónoma, especialmente basadas en
sensores avanzados y en enfoques de inteligencia artificial.

En la segunda edición (2005), el equipo de la Universidad de Stanford
logró completar con éxito la competición gracias a un sistema de
percepción más robusto basado en la fusión de sensores láser y
cámaras~\cite{thrun2006stanley}. Su principal innovación consistió en
el uso de telemetría láser para generar mapas tridimensionales del
entorno en tiempo real, demostrando que la medición precisa de
distancias ofrece una mayor fiabilidad que el análisis puramente visual
en entornos no estructurados.

El \textit{Urban Challenge} de 2007 trasladó la competición a un
entorno urbano respetando todas las normas de tráfico y compartiendo la
vía con otros vehículos, tanto autónomos como conducidos por personas.
En esta edición, los vehículos \textit{Boss} (Carnegie Mellon
University)~\cite{urmson2008boss} y \textit{Junior} (Stanford
University)~\cite{petrovskaya2009model} lograron los primeros puestos
del podio gracias al invento de David Hall: un LiDAR giratorio de 64
haces, que desempeñó un papel clave y demostró el enorme potencial de
este tipo de sensor. Este invento permitió generar nubes de puntos
tridimensionales mucho más densas y fiables, facilitando una
reconstrucción geométrica precisa del entorno en tiempo real.

Este salto en la calidad de la percepción resultó determinante, ya que
permitió pasar de sistemas que únicamente detectaban obstáculos de
forma parcial a arquitecturas capaces de comprender la escena de manera
estructurada, integrando detección, localización y planificación sobre
representaciones 3D coherentes. En este sentido, los DARPA Challenges
consolidaron la percepción como el elemento crítico en la evolución de
la conducción autónoma moderna. Estos retos no solo impulsaron el
desarrollo tecnológico, sino que establecieron una sinergia fundamental
entre los fabricantes de automóviles y el ámbito académico.

A partir de 2009, esta tecnología abandonó los laboratorios
universitarios con el nacimiento del \textit{Google Self-Driving Car
Project} (actualmente Waymo). Este hito inició una carrera comercial
sin precedentes que impulsó a la industria a estandarizar los niveles
de automatización.

\section{Niveles de automatización SAE J3016}
\label{sec:sae-j3016}

La transición desde la conducción clásica hacia la conducción
plenamente autónoma se está haciendo de forma gradual y con objeto de
establecer una clasificación clara. La Sociedad de Ingenieros de
Automoción (\textit{Society of Automotive Engineers}, SAE) estableció
en 2014 el estándar J3016, donde se definen varios niveles de
automatización, posteriormente adoptados por toda la industria,
considerando aquí su versión más reciente (SAE International,
2021)~\cite{sae2021j3016}.

Los niveles se diferencian principalmente en función de qué funciones
asume el sistema respecto al conductor humano:

\begin{itemize}
    \item \textbf{Control de movimiento:} incluye tanto el control
    lateral (dirección) como el longitudinal (aceleración y frenado)
    del vehículo.
    \item \textbf{Detección y respuesta ante objetos y eventos}
    (\textit{Object and Event Detection and Response}, OEDR):
    constituye una subtarea crítica que integra la monitorización del
    entorno mediante la detección, reconocimiento y clasificación de
    objetos, así como la ejecución de la respuesta adecuada ante los
    eventos identificados.
    \item \textbf{Planificación de maniobras:} responsable de las
    funciones tácticas, tales como la decisión de adelantar o la
    ejecución de un cambio de carril.
\end{itemize}

En la Figura~\ref{fig:niveles-automatizacion} se muestran los
diferentes niveles de automatización.

\begin{figure}[H]
    \centering
    \includegraphics[width=0.85\textwidth]{figures/niveles-automatizacion.png}
    \caption{Niveles de automatización~\cite{valle2021ecosistema}.}
    \label{fig:niveles-automatizacion}
\end{figure}

\begin{itemize}
    \item \textbf{Nivel 0:} el conductor asume el control absoluto del
    vehículo.
    \item \textbf{Nivel 1:} el sistema asiste en un único eje de
    control (longitudinal o lateral). El conductor supervisa
    permanentemente.
    \item \textbf{Nivel 2:} el sistema controla simultáneamente
    dirección y pedales, pero el conductor sigue siendo responsable del
    OEDR y debe permanecer atento. El sistema autónomo debe
    desactivarse cuando el control pase a realizarlo el conductor
    manualmente.
    \item \textbf{Nivel 3:} el sistema asume la conducción completa
    dentro de condiciones específicas, pero el conductor debe estar
    disponible para retomar el control cuando el sistema lo solicite.
    \item \textbf{Nivel 4:} el sistema conduce de forma completamente
    autónoma dentro de un ámbito operativo controlado, sin requerir
    intervención humana en su dominio de operación.
    \item \textbf{Nivel 5:} automatización plena en cualquier
    condición. El vehículo puede prescindir totalmente de controles
    manuales (sin volante ni pedales).
\end{itemize}

Desde el punto de vista de la arquitectura del sistema, esta
descomposición funcional se materializa habitualmente siguiendo el
modelo clásico de robótica móvil conocido como Percibir--Planificar--Actuar
(\textit{Sense--Plan--Act}), según el cual el vehículo se organiza en
tres módulos principales: el módulo de percepción, encargado del OEDR;
el módulo de decisión y planificación de ruta; y el módulo de control,
responsable de traducir las decisiones en comandos de dirección,
aceleración y frenado sobre los actuadores del vehículo.

El presente trabajo se centra específicamente en el módulo de
percepción, que constituye la base del flujo \textit{Sense--Plan--Act}
y el punto crítico del sistema, dado que cualquier error en esta etapa
se propaga a los módulos posteriores de planificación y control. Los
apartados siguientes describen los sensores que proporcionan la
información del entorno y los algoritmos que procesan dicha información
para detectar obstáculos y regiones transitables. Los módulos de
planificación y control quedan fuera del alcance de esta memoria.

\section{Módulo de percepción}
\label{sec:modulo-percepcion}

Este módulo es responsable de interpretar el entorno del vehículo. A
partir de un conjunto de sensores, detecta y caracteriza los elementos
estáticos y dinámicos de la escena, estimando no solo su posición y
velocidad sino también su movimiento futuro, cuya predicción resulta
especialmente compleja en el caso de agentes con comportamientos
erráticos. De forma complementaria, debe determinar la pose y velocidad
propias del vehículo respecto a ese entorno. A continuación, se
describen los tres tipos de sensores más empleados en los sistemas de
percepción actuales.

\subsection{Cámaras}
\label{sec:camaras}

Las cámaras son los sensores responsables de la \textbf{percepción
semántica} del entorno del vehículo autónomo~\cite{janai2020computer}.
Sus imágenes son esenciales para tareas como clasificación de señales
de tráfico, lectura de marcas viales e identificación de peatones. La
Figura~\ref{fig:yolo} ilustra el papel de las cámaras como sensor
responsable de la interpretación semántica de la escena.

\begin{figure}[H]
    \centering
    \includegraphics[width=0.7\textwidth]{figures/yolo-trafico.png}
    \caption{Ejemplo de salida del detector de objetos YOLO sobre una
    escena de tráfico urbano~\cite{ultralytics_collision}.}
    \label{fig:yolo}
\end{figure}

Atendiendo a su configuración geométrica, se distinguen dos modalidades
principales:

\begin{itemize}
    \item \textbf{Cámaras monoculares:} proporcionan una única
    perspectiva del entorno. No permiten estimar la distancia a los
    objetos de forma directa, aunque existen técnicas de inferencia de
    profundidad mediante redes neuronales convolucionales
    (\textit{monocular depth estimation})~\cite{ranftl2022monocular}
    cuya precisión, no obstante, sigue siendo limitada en escenarios no
    estructurados.
    \item \textbf{Cámaras estéreo:} emplean dos cámaras sincronizadas
    para calcular la profundidad triangulando a partir de la
    disparidad, definida como la diferencia de posición de un mismo
    punto entre las dos imágenes. Esta configuración permite obtener
    mediciones métricas con buena precisión a distancias cortas y
    medias; sin embargo, el error aumenta de forma cuadrática a medida
    que crece la distancia, ya que la disparidad se reduce
    progresivamente con el alcance.
\end{itemize}

Las cámaras son sensores relativamente económicos, pero requieren
sistemas de procesamiento avanzados para interpretar correctamente la
información visual. Su principal limitación es la dependencia de las
condiciones ambientales: la iluminación insuficiente o excesiva, la
niebla densa, la lluvia y la nieve degradan la calidad de las imágenes
y comprometen la fiabilidad de la percepción~\cite{bijelic2020fog}. Por
ello, es necesario complementarlos con otros sensores que permitan una
localización espacial exacta de los obstáculos.

\subsection{Radio Detection And Ranging (Radar)}
\label{sec:radar}

Funciona mediante la emisión de ondas electromagnéticas milimétricas,
habitualmente en la banda de 77~GHz en aplicaciones automotrices, que
se propagan por el espacio hasta impactar con un objeto y retornar al
sensor~\cite{patole2017radar} (véase Figura~\ref{fig:radar}).

\begin{figure}[H]
    \centering
    \includegraphics[width=0.6\textwidth]{figures/radar.png}
    \caption{Representación del funcionamiento de los sensores de
    Radar~\cite{martil_radar}.}
    \label{fig:radar}
\end{figure}

A partir del tiempo de vuelo de la señal reflejada se calcula la
distancia al obstáculo, mientras que el desplazamiento de frecuencia
entre la señal emitida y la recibida (efecto Doppler) permite estimar
de forma directa su velocidad radial. Esta capacidad de medir
simultáneamente posición y velocidad sin necesidad de procesamiento
adicional constituye una de las características distintivas del radar.

Su principal ventaja operativa es la robustez frente a condiciones
meteorológicas adversas, lo que permite mantener la detección incluso
cuando la visibilidad óptica está fuertemente
comprometida~\cite{bijelic2020fog}. Adicionalmente, destacan por su
alcance operativo, que en configuraciones de largo alcance
(\textit{Long Range Radar}, LRR) puede situarse entre los 250 y 300
metros, valores superiores a los del LiDAR comercial actual.

No obstante, el radar presenta limitaciones importantes derivadas de
la longitud de onda empleada. La principal es su baja resolución
angular, típicamente de 1--2\textdegree, lo que dificulta distinguir
objetos cercanos entre sí y determinar con precisión la forma de un
obstáculo, restringiendo su uso a la detección y seguimiento más que a
la caracterización geométrica del entorno~\cite{bilik2019radar}.
Adicionalmente, en escenarios urbanos con abundancia de superficies
metálicas (vehículos, mobiliario urbano, fachadas), las ondas
reflejadas pueden seguir trayectorias múltiples antes de retornar al
sensor (\textit{multipath}), generando detecciones fantasma o ecos
falsos cuya correcta interpretación requiere técnicas avanzadas de
filtrado y asociación de datos.

\subsection{LiDAR (Light Detection and Ranging)}
\label{sec:lidar}

Mientras que las cámaras son indispensables para la comprensión
semántica del entorno y el radar resulta clave para la detección de
objetos a larga distancia y la medición directa de la velocidad, el
LiDAR se posiciona como el sensor fundamental para la reconstrucción
geométrica precisa del entorno en tres
dimensiones~\cite{royo2019lidar}. Esta capacidad de capturar la
estructura 3D real de la escena, sin depender de las condiciones de
iluminación ni de la apariencia visual de los objetos, lo convierte en
el sensor de referencia para sistemas de percepción orientados a la
navegación autónoma segura.

\subsubsection{Principio de funcionamiento}
\label{sec:lidar-principio}

Es un sensor activo que mide distancias mediante la emisión de pulsos
láser, típicamente en longitudes de onda de 905~nm o 1550~nm. El método
estándar es el cálculo del Tiempo de Vuelo (\textit{Time of Flight},
ToF) entre la emisión del pulso y la recepción del eco reflejado por el
objeto~\cite{li2020lidar}. Conocida la velocidad de la luz
$c \approx 3 \times 10^{8}$\,m/s, la distancia $R$ al objeto se obtiene
como:
\begin{equation}
    R = \frac{c \cdot \Delta t}{2},
\end{equation}
donde $\Delta t$ es el intervalo entre la emisión del pulso y la
recepción del eco, y el factor $\tfrac{1}{2}$ refleja que la señal
recorre el camino de ida y vuelta.

Los sistemas LiDAR modernos repiten este proceso millones de veces por
segundo, generando una representación tridimensional densa del entorno
conocida como \textbf{nube de puntos}: un conjunto de puntos con
coordenadas espaciales $(X, Y, Z)$ expresadas en el sistema de
referencia local del sensor~\cite{gomes2023survey}.

Cada punto de la nube no se reduce a sus tres coordenadas espaciales,
sino que incluye información complementaria relevante para el procesado
posterior: un valor de \textbf{intensidad} que refleja la reflectividad
del material impactado, el identificador del \textbf{anillo} emisor
(\textit{ring}) que generó el haz (en sensores rotativos o en
arquitecturas \textit{solid-state} que emulan esta estructura), el
\textbf{timestamp} preciso de la adquisición, y, en sensores avanzados,
métricas adicionales como reflectividad calibrada o imagen de
infrarrojo ambiente.

Los sistemas más avanzados emplean modulación de frecuencia continua
(\textit{Frequency-Modulated Continuous Wave}, FMCW), que no solo mide
la distancia sino que permite extraer la velocidad radial del objeto
mediante el desplazamiento Doppler, de forma análoga al radar. Esta
tecnología, sin embargo, aún se encuentra en fase de desarrollo
industrial: su alto coste y la complejidad de fabricación óptica
dificultan su adopción masiva en vehículos de serie.

\subsubsection{Arquitecturas de hardware: rotativos vs.\ estado sólido}
\label{sec:lidar-hw}

Las principales arquitecturas LiDAR en uso actualmente, descritas en
detalle por Li e Ibanez-Guzman~\cite{li2020lidar}, se dividen en dos
grandes categorías: mecánicas rotativas y de estado sólido. La manera
en que el sensor barre el entorno determina su funcionalidad,
fiabilidad mecánica y coste.

\begin{itemize}
    \item \textbf{LiDAR mecánico rotativo:} emplea pares emisor-receptor
    montados sobre un eje que gira a velocidades típicas de 5--20~Hz,
    generando un barrido horizontal de 360\textdegree\
    (Figura~\ref{fig:lidar-rotativo}). Esta arquitectura, utilizada en
    sensores como el Velodyne HDL-64E o el Velodyne Alpha Prime,
    proporciona una distribución uniforme de puntos en acimut (ángulo
    horizontal medido respecto a la dirección de avance del vehículo),
    lo que la hace ideal para tareas de percepción omnidireccional. Sin
    embargo, sus partes móviles plantean desafíos de fiabilidad
    mecánica a largo plazo y son vulnerables a vibraciones. Además, su
    coste históricamente elevado y su tamaño voluminoso dificultan su
    integración estética en vehículos de producción.
\end{itemize}

\begin{figure}[H]
    \centering
    \includegraphics[width=0.55\textwidth]{figures/lidar-rotativo.png}
    \caption{Principio del LiDAR mecánico rotativo~\cite{li2020lidar}.}
    \label{fig:lidar-rotativo}
\end{figure}

\begin{itemize}
    \item \textbf{LiDAR de estado sólido} (\textit{solid-state}):
    elimina las partes móviles sustituyéndolas por tecnologías como
    \textit{Micro-Electro-Mechanical Systems} (MEMS) con micro-espejos
    oscilantes, \textit{Optical Phased Arrays} (OPA) que dirigen el haz
    electrónicamente sin elementos mecánicos, o \textit{Flash LiDAR},
    que ilumina toda la escena simultáneamente y captura el retorno con
    una matriz de detectores. Sus principales ventajas son un coste
    sustancialmente menor, mayor robustez mecánica, menor tamaño y
    posibilidad de integración embebida en parachoques y faros. Su
    limitación principal es un campo de visión más reducido que los
    rotativos (con valores típicos entre 30\textdegree\ y
    120\textdegree\ horizontales según la tecnología y el fabricante),
    lo que obliga a combinar varios sensores para conseguir cobertura
    completa.
\end{itemize}

Las arquitecturas híbridas combinan componentes de estado sólido con un
barrido mecánico simplificado. Un ejemplo destacado es la familia
Ouster OS, que rota a 10--20~Hz para lograr cobertura horizontal de
360\textdegree, posicionándose entre los LiDAR mecánicos rotativos y
los de estado sólido puros.

\subsubsection{Canales (número de haces)}
\label{sec:lidar-canales}

El número de \textbf{canales} o \textit{beams} determina la
\textbf{resolución vertical} del sensor y, en consecuencia, la densidad
de la nube de puntos generada. Esta tecnología es fundamental para la
navegación de vehículos autónomos, donde a mayor número de canales,
mayor es la riqueza de detalles y la seguridad. Los LiDAR de
\textbf{16~canales} (p.~ej.\ Velodyne VLP-16) ofrecen baja resolución y
son típicos de plataformas de bajo coste o aplicaciones poco exigentes.
Los de \textbf{64~canales} (Velodyne HDL-64E) fueron durante años el
estándar de referencia para investigación en conducción
autónoma~\cite{geiger2012kitti}. Actualmente, sensores de
\textbf{128~canales} (Ouster OS1-128 o Velodyne Alpha Prime)
proporcionan nubes mucho más densas que permiten detectar objetos
pequeños y bordillos a mayor distancia, con la contrapartida de un
coste superior, mayor ancho de banda de datos y, habitualmente, mayor
consumo energético.

\subsubsection{Características técnicas y limitaciones}
\label{sec:lidar-limitaciones}

Los parámetros principales que definen el rendimiento de un sensor
LiDAR son:

\begin{itemize}
    \item \textbf{Alcance máximo:} depende de la potencia del láser y
    de la reflectividad del objeto. En sensores de gama alta, el
    alcance puede llegar a 200--300~m sobre superficies altamente
    reflectantes, pero se reduce a 50--80~m sobre materiales oscuros o
    poco reflectantes (asfalto mojado, ropa negra).
    \item \textbf{Precisión métrica:} del orden de 2--3~cm a distancias
    medias, una de las virtudes principales del sensor frente a las
    cámaras y al radar.
    \item \textbf{Resolución angular horizontal y vertical:} los LiDAR
    de gama alta alcanzan resoluciones del orden de 0.1\textdegree,
    frente a los 1--2\textdegree\ típicos del radar.
    \item \textbf{Frecuencia de adquisición:} típicamente 10--20~Hz,
    compatible con los requisitos de tiempo real de la mayoría de
    sistemas de percepción.
    \item \textbf{Robustez ante condiciones adversas:} el rendimiento
    del LiDAR se degrada con lluvia intensa, niebla densa, nieve y
    polvo en suspensión, debido a la dispersión y absorción de la
    radiación láser por las partículas en
    suspensión~\cite{bijelic2020fog}. Esta limitación motiva la fusión
    sensorial con cámaras y radar en aplicaciones comerciales, aunque
    la robustez geométrica del LiDAR frente a variaciones de
    iluminación lo mantiene como sensor principal para la
    reconstrucción 3D en la mayoría de los \textit{pipelines} de
    investigación.
\end{itemize}

\section{Algoritmos de percepción}
\label{sec:algoritmos-percepcion}

Una vez adquirida la información del entorno a través de los sensores
descritos en la sección anterior, el vehículo autónomo debe procesarla
para construir una representación útil de la escena. Es en esta etapa
donde entran en juego los algoritmos de percepción, responsables de
transformar los datos crudos, y en particular las nubes de puntos
LiDAR, en entidades interpretables por los módulos posteriores de
planificación y control.

El módulo de percepción engloba varias subtareas, organizadas
tradicionalmente en cuatro bloques funcionales
(Figura~\ref{fig:pipeline-percepcion}): la detección de objetos
(\textit{Object Detection}), que comprende la segmentación del suelo y
la identificación de los elementos no-suelo presentes en la escena; el
reconocimiento (\textit{Object Recognition}), que asigna una categoría
semántica a cada objeto detectado; el seguimiento multi-objeto
(\textit{Multi-Object Tracking}), que mantiene la identidad y la
trayectoria de los agentes a lo largo del tiempo; y la predicción de
intención (\textit{Intention Prediction}), que estima su movimiento
futuro a corto y medio plazo.

\begin{figure}[H]
    \centering
    \includegraphics[width=0.95\textwidth]{figures/pipeline-percepcion.png}
    \caption{Pipeline de un sistema de percepción LiDAR
    clásico~\cite{li2020lidar}.}
    \label{fig:pipeline-percepcion}
\end{figure}

El presente trabajo se centra en la primera de estas etapas, la
detección de objetos, descomponiéndola en tres subtareas consecutivas:
la segmentación del suelo, la identificación de obstáculos mediante
agrupamiento de puntos no-suelo y la extracción de cajas delimitadoras
orientadas (\textit{Oriented Bounding Boxes}, OBBs) que representan de
forma compacta a cada obstáculo. Estas tres etapas constituyen el
núcleo mínimo necesario para alimentar al módulo de planificación con
la información geométrica imprescindible para una navegación segura. A
lo largo de esta sección se revisan los principales algoritmos
empleados en cada una de ellas, tanto desde el enfoque geométrico
clásico como desde el paradigma del aprendizaje profundo.

\subsection{Paradigmas: enfoque geométrico vs.\ aprendizaje profundo}
\label{sec:paradigmas}

Los algoritmos de percepción sobre nubes de puntos LiDAR pueden
agruparse en dos grandes paradigmas con filosofías radicalmente
distintas. Cada uno responde al problema de extraer información útil de
la nube de puntos con herramientas diferentes, y presenta ventajas y
limitaciones específicas que condicionan su despliegue en aplicaciones
de conducción autónoma.

\subsubsection{Enfoque geométrico clásico}
\label{sec:enfoque-geometrico}

Explota las propiedades intrínsecas de la nube de puntos (coordenadas
espaciales, normales locales, densidad, estructura de anillos) mediante
reglas y modelos matemáticos deterministas que no requieren
entrenamiento previo. Su principal virtud es la \textbf{interpretabilidad}:
cada decisión del sistema puede rastrearse hasta un criterio geométrico
explícito (un umbral de altura, una condición de verticalidad, una
distancia a un plano), lo que facilita la depuración, la justificación
formal del comportamiento y la adaptación rápida a nuevos escenarios.
Al operar directamente sobre la geometría del entorno, el enfoque
geométrico presenta una capacidad de transferencia significativamente
superior a la del aprendizaje profundo. Su aplicación a un nuevo sensor
o entorno requiere únicamente la recalibración de un pequeño conjunto
de parámetros operativos (altura del sensor, número de anillos, rango
efectivo, umbrales de altura); en contraste con el reentrenamiento
completo, la GPU y los grandes volúmenes de datos etiquetados que
exigen los métodos basados en redes neuronales.

Como contrapartida, este enfoque depende de parámetros ajustados
manualmente, puede resultar frágil ante terrenos altamente irregulares
y, al carecer de contexto semántico, tiene dificultades para distinguir
objetos con geometría similar (por ejemplo, un peatón de un poste).

\subsubsection{Enfoque basado en aprendizaje profundo}
\label{sec:enfoque-dl}

El paradigma alternativo delega la extracción de información a redes
neuronales entrenadas sobre grandes conjuntos de datos etiquetados. En
la literatura se distinguen tres grandes familias según la
representación empleada de la nube de puntos: los métodos de
\textbf{proyección 2D}, que transforman la nube en imágenes de rango
sobre las que se aplican redes convolucionales bidimensionales
(p.\ ej.\ RangeNet++); los métodos de \textbf{discretización
volumétrica}, que dividen el espacio en vóxeles (cartesianos o
cilíndricos) y aplican convoluciones 3D (VoxelNet, Cylinder3D); y los
métodos de \textbf{procesamiento directo de puntos}, que emplean
operadores invariantes a permutaciones sobre los puntos originales
(PointNet, PointNet++). Estos métodos alcanzan precisiones muy elevadas
en el dominio de entrenamiento, incorporan de forma implícita
información semántica rica y son capaces de resolver simultáneamente
múltiples subtareas (segmentación, detección y clasificación) con una
sola arquitectura.

No obstante, presentan limitaciones importantes: requieren GPU para
inferencia en tiempo real, demandan grandes volúmenes de datos
etiquetados cuya generación es costosa, y sufren una degradación
notable cuando se aplican a sensores o entornos distintos de los
empleados durante el entrenamiento, fenómeno conocido en la literatura
como \textit{domain gap} y especialmente acusado en escenarios no
urbanos o con sensores de diferente resolución. Esta dicotomía entre
precisión en dominio conocido y robustez frente a cambios de dominio es
un compromiso ampliamente reconocido en la
literatura~\cite{gomes2023survey}.

\subsubsection{Síntesis comparativa}
\label{sec:sintesis}

La Tabla~\ref{tab:comparativa-paradigmas} resume de manera estructurada
las principales diferencias y fortalezas de ambos paradigmas.

\begin{table}[H]
    \centering
    \caption{Síntesis comparativa de ambos enfoques.}
    \label{tab:comparativa-paradigmas}
    \begin{tabularx}{\textwidth}{@{}lXX@{}}
        \toprule
        \textbf{Característica} & \textbf{Enfoque geométrico} & \textbf{Aprendizaje profundo} \\
        \midrule
        Necesidad de entrenamiento & No (basado en reglas) & Sí (\textit{datasets} etiquetados) \\
        Transferencia a nuevo sensor & Recalibración de parámetros (minutos) & Reentrenamiento completo (horas/días + datos) \\
        Interpretabilidad & Total (criterio físico) & Baja (caja negra) \\
        Requisitos \textit{hardware} & Bajos (CPU eficiente) & Altos (GPU dedicada) \\
        Riqueza semántica & Limitada & Alta \\
        Robustez en dominios no vistos & Elevada (tras recalibración) & Limitada \\
        \bottomrule
    \end{tabularx}
\end{table}

\subsection{Segmentación del suelo}
\label{sec:segmentacion-suelo}

La segmentación del suelo consiste en separar los puntos de la nube
LiDAR que pertenecen a la superficie transitable de aquellos que
corresponden a obstáculos u otros elementos del entorno. Esta tarea se
ha estudiado exhaustivamente en la literatura por su papel crítico en
el \textit{pipeline} de conducción autónoma: tal como describen Gomes
et al.\ en su \textit{survey}~\cite{gomes2023survey}, el flujo de
procesamiento estándar de una nube LiDAR comprende
\enquote{una etapa de filtrado del suelo que aísla los puntos de suelo
de los no-suelo, reduciendo el volumen de datos para las etapas
posteriores de detección y clasificación de objetos, además de mejorar
la navegación y cumplir con los requisitos de tiempo real}. Sin una
segmentación fiable del suelo, cualquier algoritmo de agrupamiento
posterior mezclaría puntos de la calzada con los de vehículos y
peatones, invalidando la detección.

La literatura clasifica los algoritmos de segmentación del suelo en
cinco grandes familias, ilustradas en la
Figura~\ref{fig:seg-suelo-familias}: métodos basados en
\textbf{rejillas 2.5D}, que proyectan la nube sobre mapas de elevación
o cuadrículas de ocupación; \textbf{modelado de superficies}
(\textit{ground modelling}), que ajusta modelos matemáticos explícitos
(planos, líneas o regresiones gaussianas) a los puntos del suelo;
\textbf{análisis de puntos adyacentes y características locales}, que
explota relaciones geométricas entre puntos vecinos mediante técnicas
como \textit{channel-based}, \textit{range images}, \textit{clustering}
o \textit{region-growing}; \textbf{inferencia de orden superior},
basada en modelos probabilísticos como \textit{Markov Random Fields}
(MRF) o \textit{Conditional Random Fields} (CRF); y métodos basados en
\textbf{aprendizaje profundo}, que emplean redes convolucionales para
clasificar cada punto como suelo o no-suelo.

\begin{figure}[H]
    \centering
    \includegraphics[width=0.85\textwidth]{figures/seg-suelo-familias.png}
    \caption{Algoritmos de segmentación del suelo.}
    \label{fig:seg-suelo-familias}
\end{figure}

De estas cinco familias, el ajuste de planos (\textit{plane fitting})
presenta un conjunto de propiedades particularmente relevantes para
aplicaciones de conducción autónoma: madurez y amplia adopción en
\textit{pipelines} de producción, compatibilidad con requisitos de
tiempo real sobre CPU, y capacidad de transferencia entre sensores
mediante recalibración mínima de parámetros. Por estas razones, es el
enfoque adoptado en el presente trabajo. Las restantes familias, aunque
también válidas, presentan limitaciones que las hacen menos adecuadas
para el caso de uso considerado: las rejillas 2.5D pierden información
en terrenos con rasantes variables, los métodos de inferencia de orden
superior resultan computacionalmente costosos para \textit{pipelines}
de tiempo real, y los basados en aprendizaje profundo requieren
reentrenamiento al cambiar de sensor y \textit{hardware} específico
para inferencia.

En los apartados siguientes se revisan los dos métodos más
representativos del ajuste de planos en percepción LiDAR:
\textit{Random Sample Consensus} (RANSAC), algoritmo fundamental e
históricamente de referencia, y Patchwork++, su evolución actual basada
en ajustes locales sobre un modelo de zonas concéntricas.

\subsubsection{Plane Fitting clásico: RANSAC}
\label{sec:ransac}

Propuesto originalmente por Fischler y Bolles~\cite{fischler1981ransac}
en el contexto de la visión por computador para el ajuste de modelos en
presencia de \textit{outliers}, es una técnica iterativa genérica de
estimación robusta de parámetros de un modelo matemático. Su aplicación
a la segmentación del suelo en nubes de puntos LiDAR consiste en
ajustar, de forma iterativa, el plano tridimensional que mejor
representa la superficie transitable.

El procedimiento básico consiste en seleccionar aleatoriamente tres
puntos de la nube (el mínimo necesario para definir un plano), ajustar
el plano determinado por ellos y contar cuántos puntos del resto de la
nube se encuentran a una distancia ortogonal inferior a un umbral
prefijado, considerándolos \textit{inliers}.

\begin{figure}[H]
    \centering
    \includegraphics[width=0.55\textwidth]{figures/ransac-plane.png}
    \caption{Representación visual de una clasificación por distancia
    ortogonal~\cite{gomes2023survey}.}
    \label{fig:ransac-plane}
\end{figure}

Este proceso se repite un número suficiente de iteraciones para
garantizar con alta probabilidad que al menos una selección de tres
puntos pertenezca realmente al suelo; el plano con mayor número de
\textit{inliers} se retiene como estimación final. La naturaleza del
criterio de consenso (decidir por mayoría de \textit{inliers}) hace que
RANSAC sea robusto frente a \textit{outliers}.

No obstante, la aplicación de RANSAC a la segmentación del suelo en
entornos reales presenta limitaciones importantes. Como señalan los
autores del \textit{survey}, \enquote{el suelo no es siempre plano, por
lo que la hipótesis de un único plano puede llevar a clasificaciones
incorrectas. Además, datos ruidosos o complejos pueden degradar el
rendimiento de RANSAC}~\cite{gomes2023survey}. Por este motivo, en
entornos reales de conducción autónoma (con pendientes, baches,
bordillos o cambios de rasante) la superficie transitable no se puede
representar adecuadamente con un único plano global, lo que motiva el
desarrollo de métodos de ajuste local por regiones.

\subsubsection{Plane Fitting local: Patchwork++}
\label{sec:patchworkpp}

Para superar la limitación del plano único de RANSAC, diversos trabajos
han propuesto dividir la nube de puntos en sectores y ajustar un plano
local independiente en cada uno. El exponente más representativo de
esta línea es \textbf{Patchwork}~\cite{lim2021patchwork}, propuesto por
Lim et al.\ en 2021, que introduce una aproximación
\textit{region-wise} al ajuste de planos mediante tres contribuciones
principales:

\begin{itemize}
    \item Un \textbf{Concentric Zone Model} (CZM) que particiona el
    espacio en cuatro zonas concéntricas con tamaños angulares y
    radiales distintos por zona, adaptadas a la densidad típica de la
    nube de puntos en cada rango de distancias
    (Figura~\ref{fig:czm}). Esta partición no uniforme mitiga los
    problemas clásicos de las representaciones polares convencionales,
    donde los \textit{bins} resultan excesivamente grandes en zonas
    lejanas, lo que impide representar adecuadamente objetos pequeños
    como bordillos; mientras que los \textit{bins} próximos al sensor
    pueden contener demasiados puntos concentrados en áreas pequeñas,
    complicando el ajuste local del plano.
\end{itemize}

\begin{figure}[H]
    \centering
    \includegraphics[width=0.75\textwidth]{figures/czm.png}
    \caption{(a) Descripción mediante una malla polar uniforme.
    (b) Descripción de malla polar basada en
    CZM~\cite{lim2021patchwork}.}
    \label{fig:czm}
\end{figure}

\begin{itemize}
    \item Un método de \textbf{Region-wise Ground Plane Fitting
    (R-GPF)} que ajusta un plano local independiente en cada \textit{bin}
    mediante Análisis de Componentes Principales (PCA) sobre los $K$
    puntos de menor elevación, con velocidades al menos el doble que
    las reportadas por RANSAC.
    \item Una etapa de \textbf{Ground Likelihood Estimation (GLE)} que
    filtra los planos estimados basándose en tres características:
    \textit{uprightness} (verticalidad de la normal),
    \textit{elevation} (altura del plano) y \textit{flatness} (varianza
    de los puntos respecto al plano ajustado).
\end{itemize}

En las evaluaciones reportadas por los autores sobre el \textit{dataset}
SemanticKITTI~\cite{behley2019semantickitti} (descrito en detalle en
la Sección~\ref{sec:datasets}), Patchwork obtiene un F1 del 93.0\,\%
con una precisión del 91.38\,\% y un \textit{recall} del 93.65\,\%,
procesando \textit{frames} de aproximadamente 100.000 puntos a una
velocidad de 43.97~Hz sobre CPU.

A pesar de los buenos resultados obtenidos por Patchwork, sus propios
autores identifican limitaciones residuales que motivan el desarrollo
de una versión mejorada. Por un lado, el método presenta
sub-segmentación parcial en \textit{bins} con terrenos irregulares o
con estructuras verticales bajas cercanas al suelo (muros bajos, setos,
vallas al ras), cuyo ajuste PCA resulta contaminado por los puntos
verticales. Por otro, los umbrales fijos de la \textit{Ground
Likelihood Estimation} requieren reajuste manual entre entornos
(urbano, autopista, rural) y no aprovechan la coherencia temporal
disponible en secuencias LiDAR reales. Adicionalmente, el algoritmo no
contempla el ruido por reflexiones múltiples por debajo del suelo,
habitual en superficies altamente reflectantes (nieve, asfalto mojado),
lo que puede sesgar la estimación del plano. Para resolver estas
limitaciones, proponen en 2022 \textbf{Patchwork++}~\cite{lee2022patchworkpp},
una evolución del método original que introduce cuatro módulos
adicionales (Figura~\ref{fig:patchworkpp}):

\begin{itemize}
    \item \textbf{Reflected Noise Removal (RNR)}, que elimina los
    puntos espurios generados por reflexiones múltiples por debajo del
    suelo, problema habitual en superficies altamente reflectantes como
    la nieve o el asfalto húmedo.
    \item \textbf{Region-wise Vertical Plane Fitting (R-VPF)}, que
    detecta y filtra los planos verticales próximos al suelo antes de
    ejecutar R-GPF, evitando que estas estructuras contaminen la
    estimación del plano.
    \item \textbf{Adaptive Ground Likelihood Estimation (A-GLE)}, que
    sustituye los umbrales fijos de la GLE original por umbrales
    adaptativos calculados estadísticamente a partir de los planos
    correctamente clasificados en \textit{frames} anteriores,
    permitiendo al algoritmo adaptarse automáticamente a distintos
    entornos (urbano, autopista, rural) sin reajuste manual de
    parámetros.
    \item \textbf{Temporal Ground Revert (TGR)}, etapa final que
    re-examina los \textit{bins} inicialmente clasificados como no-suelo
    y los reclasifica como suelo si sus características son consistentes
    con las del entorno medio del \textit{frame}.
\end{itemize}

\begin{figure}[H]
    \centering
    \includegraphics[width=0.85\textwidth]{figures/patchworkpp.png}
    \caption{Diagramas de flujo de Patchwork y
    Patchwork++~\cite{lee2022patchworkpp}.}
    \label{fig:patchworkpp}
\end{figure}

La combinación de estas mejoras permite a Patchwork++ alcanzar un F1
del 96.51\,\% sobre SemanticKITTI, frente al 95.58\,\% reportado por
Patchwork, con una mejora especialmente notable en el \textit{recall}
(+1.4 puntos). Por su combinación de precisión, velocidad sobre CPU y
robustez frente a distintos entornos, Patchwork++ constituye
actualmente el estado del arte en segmentación del suelo basada en
métodos geométricos, y es el algoritmo adoptado como base del
\textit{pipeline} propuesto en el presente trabajo.

A pesar de las mejoras introducidas por Patchwork++, el método mantiene
limitaciones intrínsecas al paradigma de ajuste de planos por regiones.
El módulo TGR, aunque mejora notablemente el \textit{recall}, puede
introducir falsos positivos al recuperar como suelo \textit{bins} con
baja varianza vertical pero no pertenecientes al suelo real, con un
\textit{trade-off} que los autores consideran aceptable. Además, en
trabajos del mismo grupo~\cite{oh2022travel} se señala que
\enquote{todos los métodos basados en segmentación del suelo, incluido
Patchwork, se han centrado únicamente en la separación
\textit{ground/non-ground} sin examinar la transitabilidad del terreno,
es decir, la capacidad de identificar un área sobre la que un vehículo
puede desplazarse con seguridad}, lo cual implica que elementos como
bordillos o pequeños obstáculos en la calzada pueden quedar absorbidos
dentro del modelo del plano al cumplir los criterios de
\textit{uprightness}, \textit{elevation} y \textit{flatness}.

La detección explícita de estos elementos requiere técnicas
complementarias específicas, que se revisan en el siguiente apartado.

\subsection{Detección de bordillos}
\label{sec:deteccion-bordillos}

La detección de bordillos constituye una tarea de percepción específica
y especialmente relevante para la navegación autónoma, al situarse en
la frontera entre el suelo transitable y los elementos no navegables
del entorno. Los bordillos delimitan físicamente la calzada respecto a
aceras, medianas y zonas peatonales, por lo que su correcta
identificación resulta imprescindible para que el vehículo respete los
límites de la vía y planifique trayectorias seguras, incluso en
ausencia de marcas viales pintadas.

Sin embargo, su detección entraña una dificultad notable debido a su
geometría: con alturas típicas comprendidas entre 10 y 20~cm, los
bordillos se sitúan en el límite de resolución métrica de los sensores
LiDAR convencionales y, desde un punto de vista puramente geométrico,
son indistinguibles de un suelo localmente irregular. Esta limitación
ha impulsado el desarrollo de técnicas específicas en la literatura,
cuya evolución histórica ha sido sintetizada por Romero et
al.~\cite{romero2021curb} en su revisión del campo.

La revisión realizada por dichos autores incluye contribuciones tanto
sobre sistemas de percepción en tiempo real para vehículos autónomos
(con sensores LiDAR rotativos a 10--20~Hz) como sobre sistemas de
\textit{Mobile Laser Scanning} (MLS), configuraciones que combinan
LiDAR con GNSS-RTK e IMU para la generación de cartografía 3D densa en
post-procesado. Aunque las condiciones operativas de ambos entornos
difieren (tiempo real vs.\ \textit{offline}, nube de puntos automotriz
de densidad variable vs.\ nube MLS densa y sistemática), numerosos
principios metodológicos desarrollados en el contexto MLS son
conceptualmente transferibles al paradigma automotriz. Los siguientes
apartados revisan las principales estrategias propuestas, organizadas
según el principio metodológico subyacente.

\subsubsection{Métodos basados en características geométricas radiales}
\label{sec:curbs-radiales}

Esta estrategia es la más extendida en sistemas de tiempo real por su
bajo coste computacional. Se basa en explotar la estructura radial
intrínseca de los sensores LiDAR rotativos, analizando los puntos de
forma secuencial según el anillo (\textit{ring}) al que pertenecen.
Dentro de este paradigma conviven dos variantes complementarias, que
comparten la misma intuición: al barrer un bordillo, el sensor registra
un \textbf{salto vertical} entre puntos contiguos (ya sea entre
\textit{rings} o dentro del mismo \textit{ring}) que supera un umbral
característico, permitiendo identificar la discontinuidad.

Los métodos \textbf{inter-anillo} comparan métricas de altura entre
puntos de anillos adyacentes $(P_i, P_{i+1})$ para un mismo acimut,
detectando la discontinuidad cuando el anillo interior cae sobre la
calzada y el exterior sobre la acera. Los métodos \textbf{intra-anillo},
en cambio, analizan las transiciones entre puntos consecutivos dentro
de un mismo anillo por separado, interpretando el bordillo como una
ruptura local en la secuencia ordenada de puntos del barrido láser.

El trabajo impulsor de esta línea es el de Tan et
al.~\cite{tan2014curb}, que formalizan la detección de bordillos sobre
un Velodyne HDL-64E combinando el análisis inter-anillo con información
de cámara a color. Sobre la nube LiDAR, el método procesa pares de
puntos de anillos consecutivos $(P_i, P_{i+1})$ correspondientes a una
misma dirección angular y calcula el salto de altura entre ellos:
\begin{equation}
    \Delta z = \left| z_{i+1} - z_i \right|.
\end{equation}
Cuando el sensor barre un bordillo, el anillo interior $i$ cae sobre la
calzada ($z \approx 0$) mientras el exterior $i+1$ cae sobre la acera
($z \approx +15$\,cm), de modo que $\Delta z$ supera un umbral
$H_{\min}$ característico y el punto se marca como candidato a bordillo.
Posteriormente, se valida con una máscara semántica extraída de la
imagen, lo que reduce falsos positivos en zonas con vegetación o
sombras.

La fusión LiDAR--cámara aporta robustez frente a geometrías ambiguas,
pero introduce también sus principales limitaciones: exige una
calibración espacio-temporal precisa entre ambos sensores, se degrada
con niebla o baja iluminación, y restringe el alcance efectivo de la
detección a aproximadamente 30~m, por debajo del rango útil real del
LiDAR aislado.

Chen et al.~\cite{chen2015velodyne} extienden el alcance de la
detección hasta 50~m trabajando exclusivamente con el Velodyne HDL-64E,
renunciando a la información de cámara y adoptando un análisis
intra-anillo puro: procesan cada anillo $P_m$ como una secuencia de
puntos consecutivos. Su \textit{pipeline} combina tres fases:
(i)~extracción de \textit{feature points} por anillo mediante tres
descriptores locales: la \textit{smoothness}, que mide la variación del
punto respecto a sus vecinos del mismo anillo; la \textit{smooth arc
length}, que cuantifica la longitud de los arcos suaves a izquierda y
derecha sin cambios bruscos de geometría; y la \textit{height difference},
definida como
\begin{equation}
    \Delta z = Z_{\max} - Z_{\min}, \quad Z_{\max}, Z_{\min} \in S \subset P_m,
\end{equation}
donde $S$ es el conjunto de puntos consecutivos del mismo anillo
alrededor de $p_{m,i}$;
(ii)~selección de semillas iniciales mediante criterio de distancia y
Transformada de Hough sobre una rejilla polar;
(iii)~modelado iterativo del bordillo mediante \textit{Gaussian Process
Regression} (GPR), que permite representar tanto tramos rectos como
curvas de forma continua. La GPR iterativa maneja bordillos curvos con
precisión, algo que los métodos basados en Hough no capturan; las tres
características locales combinadas son más robustas que un único
umbral de altura, y la regresión GPR permite recuperar tramos
parcialmente ocluidos gracias a su capacidad predictiva.

Como contrapartida, el método alcanza una latencia de
$\sim$69~ms por \textit{frame} sobre CPU i5 ($\sim$14~Hz), claramente
superior a la de los métodos inter-anillo más simples; la calibración
de los cinco hiperparámetros de la GPR ($l$, $\sigma_f^2$,
$\sigma_n^2$, $T_{\text{model}}$, $T_{\text{data}}$) requiere un
proceso previo de aprendizaje sobre curvas etiquetadas manualmente; y,
además, el \textit{pipeline} presupone que el vehículo circula siempre
con bordillos a ambos lados, simplificación que deja de cumplirse en
autopistas o zonas rurales.

Zhang et al.~\cite{yang2013mls} proponen el \textit{pipeline} más
completo de esta familia, orientado específicamente a conducción
autónoma, validado sobre un Velodyne HDL-32E y estructurado en dos
fases diferenciadas. La primera fase consiste en una segmentación
previa de la calzada mediante un \textit{sliding-beam model} capaz de
reconocer el trazado incluso en intersecciones complejas (en forma de
T, + e Y), donde los bordillos dejan de ser continuos. La segunda fase
aplica una detección de bordillos por segmento basada en cuatro
criterios geométricos intra-anillo calculados sobre cada \textit{ring}
$l$ del sensor: continuidad horizontal $\delta_{xy,l}$, continuidad
vertical $\delta_{z,l}$, ángulo entre vectores consecutivos
$\theta_{i,l}$ y altura mínima del salto $H_c$. La continuidad
horizontal se deriva físicamente de la geometría del sensor como
\begin{equation}
    \delta_{xy,l} = H_s \cdot \cot \theta_{f,l} \cdot \frac{\pi \, \theta_a}{180},
\end{equation}
donde $H_s$ es la altura de montaje del LiDAR, $\theta_{f,l}$ el ángulo
vertical del anillo $l$ y $\theta_a$ la resolución angular horizontal;
esta magnitud representa la distancia horizontal esperada entre dos
puntos consecutivos del mismo anillo cuando ambos están sobre la
calzada. La continuidad vertical se obtiene de manera análoga como
\begin{equation}
    \delta_{z,l} = \delta_{xy,l} \cdot \sin \theta_{f,l},
\end{equation}
que representa el desplazamiento vertical esperado entre los mismos
puntos sobre un suelo plano. Si la diferencia real de altura
$\left| z_i - z_{i-1} \right|$ entre dos puntos consecutivos del anillo
$l$ supera $\delta_{z,l}$, se identifica una discontinuidad intra-anillo
compatible con un bordillo.

El hecho de que todos los umbrales se deriven explícitamente de
parámetros físicos del sensor, en lugar de ajustarse empíricamente,
permite una transferibilidad sensor-a-sensor significativamente mayor
que en las aproximaciones anteriores. Entre sus limitaciones se
encuentran la necesidad de información adicional DGPS + INS para
compensar el \textit{roll} y \textit{pitch} del vehículo ---añadiendo
dependencia de sensores externos al LiDAR---, la asunción inicial de
calzada plana en el filtrado \textit{on-road} (que puede fallar en
terrenos con pendiente acusada o rasantes irregulares) y la presencia
de tres hiperparámetros del \textit{sliding-beam model} ($D_b$, $d_b$,
$\theta_r$) cuya calibración requiere una búsqueda empírica extensiva
---hasta 60 experimentos por escenario en el \textit{paper} original---
para encontrar el ajuste óptimo.

En resumen, los tres métodos alcanzan latencias compatibles con
requisitos de tiempo real sobre CPU sin necesidad de GPU, no requieren
entrenamiento supervisado ni \textit{datasets} etiquetados masivos, y
producen decisiones directamente trazables a umbrales geométricos
explícitos (altura, ángulo, continuidad), lo que los convierte en
soluciones interpretables y fácilmente depurables. Como contrapartida,
dependen de la indexabilidad de los puntos por anillo: cualquier cambio
en el número de canales del sensor (16, 32, 64 o 128) requiere
recalibrar los umbrales, y la aplicación a LiDAR de estado sólido puros
(\textit{Flash}, OPA) exige reconstruir la estructura de líneas
equivalente a partir del patrón de escaneo específico del sensor, tarea
que no siempre resulta trivial. Además, son frágiles frente a la
variación de densidad con la distancia: los anillos se separan
progresivamente en el espacio a medida que aumenta el rango, con lo que
el ruido relativo en los umbrales de continuidad crece y el rendimiento
se degrada en las zonas lejanas al sensor. Por último, presuponen una
continuidad geométrica relativamente bien definida del bordillo:
funcionan adecuadamente en escenarios urbanos estructurados con
bordillos largos y rectos o suavemente curvados, pero su rendimiento
cae en entornos \textit{off-road}, en tramos con bordillos muy bajos
(inferiores a 10~cm, próximos al límite de resolución métrica del
sensor) o cuando el salto vertical queda diluido por la presencia de
vegetación baja, arena o nieve en el borde de la calzada.

Cabe mencionar a Yang et al.~\cite{yang2013mls}, quienes extienden el
paradigma de análisis por líneas a datos MLS más densos (Optech Lynx
Mobile Mapper), particionando la nube en secciones transversales de la
calzada mediante discontinuidades del tiempo GPS de los pulsos láser.
Sobre cada sección aplican un operador de ventana móvil análogo al
análisis intra-anillo, pero introducen dos aportaciones distintivas
respecto a los métodos revisados. La primera es el uso del cambio de
\textbf{densidad de puntos} como \textit{feature} discriminante
adicional: la cara vertical del bordillo recibe más pulsos por unidad
angular que la superficie plana adyacente, lo que produce una
concentración anómala de puntos detectable aunque el salto de altura se
aproxime al umbral de ruido. La segunda es el uso del cambio acumulado
de pendiente $\Delta \theta$ para distinguir tres tipos geométricos de
bordillo (vertical, inclinado y redondeado), frente al criterio binario
bordillo-o-no-bordillo de los demás métodos. Finalmente, una etapa de
\textit{tracking} determinista y refinamiento consolida los puntos en
entidades 3D continuas. Aunque el método fue diseñado para datos MLS de
alta densidad y post-procesamiento \textit{offline}, sus tres
aportaciones ---densidad como \textit{feature}, clasificación por tipos
de bordillo y \textit{tracking} multi-\textit{frame}--- son
conceptualmente transferibles al LiDAR automotriz en tiempo real,
siempre que se sustituya la partición por tiempo GPS por una
indexación por \textit{ring} del sensor.

\subsubsection{Análisis estadístico de la superficie (normales y PCA)}
\label{sec:curbs-pca}

A diferencia de la estrategia anterior, esta familia es agnóstica
respecto a la arquitectura del sensor: trata la nube como una colección
desestructurada de puntos en el espacio $(X, Y, Z)$ y analiza la
topología local de la superficie.

Mediante PCA aplicado sobre vecindarios locales (bolas esféricas o
$k$-\textit{nearest neighbors}) de cada punto se extraen dos magnitudes
clave: la orientación del vector normal $\vec{n}$ al plano de ajuste
local y la curvatura local
\begin{equation}
    c = \frac{\lambda_3}{\lambda_1 + \lambda_2 + \lambda_3},
\end{equation}
que cuantifica el grado de planitud del entorno. En una calzada plana,
las normales apuntan aproximadamente en dirección vertical; sobre la
cara lateral de un bordillo, la normal gira bruscamente hacia la
horizontalidad y la curvatura se dispara, revelando la discontinuidad
sin necesidad de información sobre la estructura radial del sensor.

Hervieu \& Soheilian abordan la detección y reconstrucción de bordillos
sobre datos de \textit{Mobile Laser Scanning} (sistema Stereopolis del
IGN) empleando un enfoque basado en la orientación de las normales
locales. Para cada punto de la nube, calculan el vector normal mediante
ajuste de plano por mínimos cuadrados sobre un vecindario esférico de
radio $R$. La característica discriminante del método es la distancia
angular a la normal del suelo:
\begin{equation}
    \theta = \arccos \left( V_p \cdot \vec{z} \right),
\end{equation}
donde $\vec{z}$ es la dirección vertical. En superficie plana de
calzada $\theta \approx 0^{\circ}$, mientras que en la cara lateral de
un bordillo $\theta$ crece hasta aproximadamente 65--88\textdegree. Para
paliar la inestabilidad del PCA con pocos vecinos, el método adapta
dinámicamente el radio $R_p$ a la densidad local de cada punto.
Posteriormente, un modelo predicción-observación basado en filtro de
Kalman propaga la posición estimada del bordillo a lo largo de la
calzada mediante un ajuste polinómico de grado $D$, lo que permite
reconstruir tramos ocluidos por obstáculos extrapolando el modelo hasta
que se recupera una observación válida. Entre sus ventajas destaca ser
uno de los pocos métodos que aborda explícitamente oclusiones. Sus
principales limitaciones residen en la dependencia de un radio $R$
inicial adecuado para el cálculo de normales (demasiado pequeño produce
normales ruidosas; demasiado grande pierde detalle cerca del sensor),
en el coste computacional del ajuste de planos punto a punto sobre
nubes MLS densas (no compatible con tiempo real automotriz), y en la
configuración empírica de múltiples parámetros del modelo Kalman
($\delta$, $\epsilon$, $L_M$, \textit{gains}, etc.) que requieren
ajuste por escenario.

Sus ideas fundamentales son directamente aplicables al LiDAR automotriz
en tiempo real: el cálculo de normales locales por PCA sobre
vecindarios esféricos constituye una herramienta estándar,
implementable eficientemente con estructuras KD-\textit{tree}; el uso
del ángulo con respecto a la vertical como \textit{feature}
discriminante es independiente del tipo de sensor; y el modelo Kalman
de predicción-observación es compatible con los requisitos temporales
de la percepción automotriz. Las principales adaptaciones necesarias
son la reducción previa de la densidad mediante voxelización (para
acotar el coste del PCA punto a punto) y el ajuste del radio $R$ a la
densidad real del LiDAR automotriz, sustancialmente menor a la del MLS.
Cabe destacar que el presente trabajo adapta precisamente este tipo de
análisis PCA al contexto del LiDAR automotriz para validar los
bordillos detectados en la etapa geométrica previa.

Más recientemente, Zhao et al.\ combinan el análisis PCA con un
criterio multiescala de dimensionalidad (\textit{Multi-scale
Dimensionality Criterion}, MSDC), calculando razones de valores propios
$p_i = \lambda_i / (\lambda_1 + \lambda_2 + \lambda_3)$ a distintas
escalas de vecindario para clasificar previamente los puntos en
categorías 1D (líneas, como ramas de árbol), 2D (superficies, como
calzada) y 3D (volúmenes, como vegetación densa). Esta clasificación
permite separar explícitamente la vegetación del suelo antes de aplicar
la detección de bordillos, abordando uno de los principales obstáculos
de los métodos anteriores en entornos urbanos con arbolado. Sobre los
puntos clasificados como suelo, combinan tres características
complementarias ---intensidad de reflectancia filtrada con el método de
Otsu, gradiente de elevación sobre un modelo digital de superficie
(DEM) y variación de pendiente--- y refinan los bordes detectados
mediante ajuste por \textit{splines} cuadráticas de Bézier, reportando
tasas de completitud y corrección superiores al 98\,\% sobre bordes
izquierdo y derecho en entornos urbanos complejos con vegetación,
oclusiones y curvas.

Las principales ventajas del método son su robustez frente a vegetación
gracias al MSDC, la combinación de información geométrica y fotométrica
---que aporta discriminación adicional---, el manejo explícito de
bordillos curvos mediante el ajuste final por Bézier y la
disponibilidad \textit{open access} del trabajo. Como contrapartida, el
método fue diseñado y validado sobre MLS de alta densidad (RIEGL
VUX-1HA), no sobre LiDAR automotriz de baja densidad; el cálculo
multi-escala de PCA multiplica el coste computacional por el número de
escalas evaluadas, alejándose del tiempo real sobre CPU embebido;
depende de la calibración de intensidad del sensor ---que varía entre
modelos y degrada con lluvia o suciedad---; requiere semi-supervisión
para ajustar la separabilidad entre clases en el MSDC; y el ajuste por
Bézier asume bordillos suaves y sin discontinuidades bruscas, por lo
que falla en intersecciones complejas.

\subsubsection{Métodos basados en extracción estructural paramétrica}
\label{sec:curbs-estructural}

Una tercera familia de métodos no se centra en la identificación de
puntos individuales de bordillo, sino en la detección de estructuras
geométricas coherentes (líneas, curvas paramétricas, planos) que
representen el borde de la vía de forma continua. El fundamento es que
los bordillos raramente son puntos aislados: típicamente forman líneas
o curvas suaves que delimitan la calzada, y por tanto la detección
puede beneficiarse de imponer esta restricción estructural.

Las herramientas matemáticas más empleadas son la \textbf{Transformada
de Hough}, que identifica alineaciones de puntos en espacios
paramétricos, y los algoritmos de \textbf{crecimiento de regiones}
(\textit{region growing}), que agrupan puntos candidatos en función de
criterios de continuidad. Fernández et al.\ (2013) aplicaron la
Transformada de Hough para garantizar que las detecciones de bordillos
formen líneas continuas, eliminando detecciones espurias que no
presenten coherencia estructural con sus vecinas. Kumar et al.\ (2013)
propusieron una combinación de gradientes y \textit{region growing}
para extraer bordes de carretera, mientras que Zai et al.\ (2017)
emplearon supervóxeles y \textit{graph cuts} para segmentar bordillos
en nubes de puntos densas procedentes de \textit{Mobile Laser Scanning}.

La principal virtud de este enfoque es su robustez frente a ruido
puntual: al exigir que las detecciones cumplan una restricción
estructural global, los falsos positivos aislados son descartados. Sin
embargo, estos métodos asumen que el entorno es estructurado y regular,
típicamente entornos urbanos con calles rectas y bordillos paralelos.
En escenarios \textit{off-road} o en entornos no estructurados (como
los del \textit{dataset} GOOSE), donde los bordes son irregulares,
discontinuos o no responden a formas paramétricas simples, estos
métodos presentan un rendimiento limitado.

Rodríguez-Cuenca et al.~\cite{ericson2005collision} proponen una
aproximación alternativa basada en la rasterización de la nube MLS
sobre una rejilla 2D de $5 \times 5$\,cm. Para cada celda calculan dos
magnitudes complementarias: un mapa de \textbf{diferencia de alturas}
(nDSM, \textit{normalized Digital Surface Model}) que aproxima la
altura del obstáculo sobre el suelo, y un mapa de \textbf{densidad de
puntos} que refleja la orientación de las superficies respecto al haz
láser. Mediante umbralización combinada y procesamiento morfológico se
extraen píxeles candidatos a bordillo; una búsqueda posterior de
\textbf{elementos lineales} descarta agrupaciones sin forma coherente.
La detección de bordes se realiza rotando la nube $\pm 45^{\circ}$
sobre el eje X, de modo que las caras verticales y horizontales del
bordillo queden inclinadas simétricamente y los bordes se identifiquen
como \textbf{cambios de signo de pendiente} entre puntos vecinos.
Finalmente, incorporan una etapa de estimación de bordillos ocluidos
para reconstruir tramos parcialmente tapados por vehículos o mobiliario
urbano.

\begin{itemize}
    \item \textbf{Ventajas:} tasas de detección superiores al 95\,\% en
    dos \textit{datasets} urbanos diferentes; maneja explícitamente
    bordillos rectos y curvos; detecta tanto el borde superior como el
    inferior; recupera bordillos ocluidos por obstáculos dinámicos; no
    requiere cálculo costoso de normales 3D al operar sobre imágenes
    2D.
    \item \textbf{Limitaciones:} orientado exclusivamente a
    post-procesamiento de MLS, no a tiempo real sobre vehículo en
    marcha; depende de umbrales ajustados empíricamente ($H_{\min}$,
    $H_{\max}$, densidad mínima, ancho lineal\ldots) específicos por
    \textit{dataset}; la rasterización a $5 \times 5$\,cm asume nubes
    muy densas, inviable con la densidad típica de un LiDAR automotriz;
    no validado en escenarios \textit{off-road} ni en condiciones no
    estructuradas; la rotación de $\pm 45^{\circ}$ asume que el suelo
    es aproximadamente plano, por lo que degrada en pendientes
    acusadas.
\end{itemize}

Finalmente, los métodos más recientes combinan información de múltiples
sensores y emplean técnicas de aprendizaje profundo para abordar los
casos en los que el LiDAR aislado resulta insuficiente: bordillos muy
bajos cercanos al umbral de ruido, superficies oscuras o de baja
reflectividad, y condiciones meteorológicas adversas. El principio de
la fusión sensorial consiste en complementar la geometría proporcionada
por el LiDAR con la información semántica y cromática de las cámaras
RGB: el LiDAR aporta el salto geométrico vertical, mientras que la
cámara confirma la presencia del bordillo mediante diferencias de
textura, color o la continuidad de marcas viales asociadas.

En el ámbito del aprendizaje profundo, la publicación del
\textit{dataset} 3D-Curb~\cite{zhao2024curbnet} ---con anotaciones
específicas de bordillos sobre nubes de puntos--- ha permitido el
entrenamiento de redes neuronales capaces de identificar bordillos de
forma implícita a partir de ejemplos etiquetados. Las principales
limitaciones de esta familia son la complejidad de implementación
(calibración espacio-temporal entre sensores heterogéneos, gestión de
tasas de muestreo distintas), la necesidad de \textit{hardware} de
procesamiento gráfico (GPU) para los métodos basados en redes, y la
dependencia de \textit{datasets} etiquetados específicos que cubran el
dominio de despliegue. Esta dependencia de datos supervisados es
especialmente crítica para bordillos, ya que el etiquetado de puntos
individuales en nubes LiDAR es laborioso y los \textit{datasets}
disponibles (3D-Curb sobre SemanticKITTI) cubren principalmente
entornos urbanos estructurados, lo que compromete la generalización a
escenarios no vistos.

\subsection{Identificación de obstáculos mediante agrupamiento}
\label{sec:clustering}

Una vez segmentada la superficie, los puntos restantes en la nube se
consideran candidatos a obstáculos. Sin embargo, para que futuros
sistemas de planificación de trayectorias puedan ejecutar maniobras
evasivas o de seguimiento, es necesaria una etapa de agrupamiento
(\textit{clustering}) que consolide los puntos no-suelo en clústeres
coherentes, cada uno correspondiente a un obstáculo físico distinto.

Aunque existen métodos clásicos como el \textbf{Euclidean Clustering}
(basado en distancias euclídeas simples y crecimiento regional) o el
\textbf{K-Means} (que exige prefijar el número de clústeres), la
naturaleza ruidosa y dispersa de los datos LiDAR en entornos reales
favorece el uso de \textbf{enfoques basados en densidad}, capaces de
identificar obstáculos de forma, tamaño y número variables sin
conocimiento previo del entorno.

\subsubsection{Density-Based Spatial Clustering of Applications with Noise (DBSCAN)}
\label{sec:dbscan}

DBSCAN~\cite{yabroudi2022dbscan} es uno de los pilares de la percepción
en robótica gracias a su capacidad de identificar clústeres de
\textbf{forma arbitraria}, sin asumir compacidad ni convexidad; su
\textbf{robustez intrínseca al ruido}, que descarta automáticamente los
retornos espurios; y por no requerir fijar \textbf{el número de
clústeres a priori}, lo cual es fundamental en entornos dinámicos donde
la cantidad de obstáculos es desconocida y cambia entre
\textit{frames}.

El funcionamiento del algoritmo se rige por la definición de una
\enquote{vecindad densa} mediante dos parámetros fundamentales:

\begin{itemize}
    \item $\epsilon$ (\textit{Epsilon}): representa el radio de
    búsqueda o distancia máxima para considerar que dos puntos son
    vecinos.
    \item $\mathit{minPts}$ (\textit{Minimum Points}): define el número
    mínimo de puntos que deben encontrarse dentro del radio $\epsilon$
    para constituir una región densa o núcleo.
\end{itemize}

Bajo estos criterios, el algoritmo clasifica cada punto de la nube en
tres categorías:

\begin{itemize}
    \item \textbf{Puntos Núcleo} (\textit{Core Points}): aquellos que
    contienen al menos $\mathit{minPts}$ en su vecindad $\epsilon$. Son
    el centro de los objetos detectados.
    \item \textbf{Puntos Frontera} (\textit{Border Points}): se
    encuentran dentro del radio $\epsilon$ de un punto núcleo, pero no
    poseen la densidad suficiente para generar un nuevo núcleo por sí
    mismos.
    \item \textbf{Ruido} (\textit{Noise Points}): puntos aislados que
    no cumplen ninguno de los criterios anteriores y son descartados
    automáticamente.
\end{itemize}

\begin{figure}[H]
    \centering
    \includegraphics[width=0.55\textwidth]{figures/dbscan-types.png}
    \caption{Tipos de puntos en DBSCAN~\cite{yabroudi2022dbscan}.}
    \label{fig:dbscan}
\end{figure}

La aplicación de DBSCAN a la detección de obstáculos en el contexto del
OEDR aporta ventajas estratégicas destacables. En primer lugar, su
capacidad de filtrar ruido de los retornos espurios del LiDAR
(provocados por partículas en suspensión, lluvia o reflejos) evita
falsas detecciones que podrían comprometer la seguridad del vehículo.
En segundo lugar, el enfoque basado en densidad permite agrupar
coherentemente estructuras con discontinuidades o formas alargadas,
como camiones articulados, vallas longitudinales o hileras de vehículos
estacionados, evitando la fragmentación que suelen presentar los
métodos de agrupamiento euclídeo tradicionales cuando la nube contiene
huecos por oclusión o baja reflectividad local.

Sin embargo, su aplicación directa sobre la nube de puntos enfrenta el
desafío de la \textbf{densidad variable con la distancia}. Debido a la
resolución angular del sensor, la densidad de puntos decrece de forma
cuadrática respecto a la distancia. Un valor de $\epsilon$ constante
resulta problemático: si es demasiado reducido, los objetos lejanos son
descartados como ruido; si es excesivamente amplio, objetos cercanos
(como dos peatones próximos) corren el riesgo de ser fusionados en un
único clúster. Esta limitación motiva el uso de variantes adaptativas o
enfoques de \enquote{multiepsilon} en sistemas de percepción avanzados.
En \textit{pipelines} de percepción modernos también es habitual
voxelizar previamente la nube para homogeneizar la densidad antes de
aplicar DBSCAN, reduciendo tanto el coste computacional como la
sensibilidad de los parámetros al rango.

\subsection{Representación de obstáculos: cajas delimitadoras orientadas (OBBs)}
\label{sec:obbs}

Tras identificar entidades discretas en la nube de puntos mediante
algoritmos de agrupamiento, es necesario abstraerlas en representaciones
geométricas compactas que puedan ser procesadas eficientemente por los
módulos posteriores. Trabajar con miles de puntos por obstáculo resulta
inviable para la planificación de trayectorias y la evitación de
colisiones, que requieren descripciones de baja dimensionalidad para
mantener latencias compatibles con niveles de automatización SAE 3 o
superiores.

La literatura de geometría computacional ofrece diversos volúmenes
delimitadores (esferas, AABB, OBB, $k$-DOP, \textit{convex hull}),
ordenables según su compromiso entre ajuste geométrico y coste
computacional~\cite{ericson2005collision}.

\begin{figure}[H]
    \centering
    \includegraphics[width=0.75\textwidth]{figures/bounding-volumes.png}
    \caption{Tipos de volúmenes envolventes~\cite{ericson2005collision}.}
    \label{fig:bounding-volumes}
\end{figure}

En conducción autónoma, las cajas delimitadoras se han consolidado como
representación dominante por su balance entre ajuste, coste y
compatibilidad con los módulos posteriores.

\subsubsection{AABB vs.\ OBB}
\label{sec:aabb-obb}

Existen dos enfoques principales para la construcción de cajas
delimitadoras. El primero, la \textbf{Axis-Aligned Bounding Box (AABB)},
es el más sencillo desde el punto de vista computacional: se obtiene a
partir de los valores mínimos y máximos de las coordenadas $(X, Y, Z)$
del clúster en el sistema de referencia del vehículo. Su coste es
prácticamente despreciable, lo que la hace especialmente atractiva en
sistemas en tiempo real. No obstante, cuando el obstáculo está
orientado de forma oblicua respecto al eje longitudinal del vehículo,
la AABB introduce un volumen significativo de espacio vacío en sus
esquinas. Esta sobreestimación degrada la precisión geométrica y puede
afectar negativamente a la eficiencia del
planificador~\cite{ericson2005collision}.

Para superar esta limitación, se emplea la \textbf{Oriented Bounding
Box (OBB)}, que incorpora un grado de libertad adicional: la
orientación o ángulo de guiñada $\theta$. De este modo, las aristas de
la caja se alinean con la dirección principal del clúster en lugar de
con los ejes del vehículo. El resultado es un ajuste mucho más preciso
y proporciona una representación más fiel tanto del tamaño como de la
orientación del obstáculo.

Desde un punto de vista formal, una OBB puede representarse mediante el
vector de parámetros:
\begin{equation}
    \mathbf{obb} = (c_x, c_y, c_z, l, w, h, \theta),
\end{equation}
donde $(c_x, c_y, c_z)$ define el centro del volumen, $(l, w, h)$
corresponden a sus dimensiones (largo, ancho y alto), y $\theta$ es el
ángulo de guiñada respecto al eje $X$ del vehículo. Esta parametrización
compacta permite a los módulos de planificación calcular distancias de
seguridad con precisión, realizar comprobaciones de colisión mediante
tests geométricos eficientes y estimar dinámicas del obstáculo a partir
de su evolución temporal entre \textit{frames} consecutivos.

\subsubsection{Cálculo de la OBB mediante PCA}
\label{sec:obb-pca}

El método más extendido para calcular una OBB a partir de un clúster de
puntos LiDAR consiste en aplicar Análisis de Componentes Principales
(PCA) sobre las coordenadas horizontales $(X, Y)$. Este enfoque se
apoya en una hipótesis razonable en entornos de conducción: los
obstáculos presentan una estructura aproximadamente vertical respecto
al suelo, por lo que su orientación puede capturarse adecuadamente en
el plano 2D. Implementado de forma generalizada en librerías como
PCL~\cite{rusu2011pcl}, este método destaca por su bajo coste
computacional y por ofrecer resultados satisfactorios para la mayoría
de objetos convexos en entornos urbanos.

Sin embargo, presenta una limitación importante: su sensibilidad a
\textit{outliers}, ya que puntos aislados alejados del centroide pueden
sesgar significativamente la orientación estimada. Por ello, en
aplicaciones industriales suele combinarse con etapas previas de
filtrado (como mediana o recorte por percentiles) o sustituirse por
métodos más robustos, como el algoritmo de \textit{rotating
calipers}~\cite{toussaint_calipers} aplicado sobre el \textit{convex
hull} 2D del clúster, que garantiza la obtención de la caja de área
mínima a costa de un ligero incremento en el coste computacional.

Tras el análisis de las estrategias algorítmicas para la segmentación y
representación de obstáculos, resulta necesario establecer un marco de
evaluación que permita medir su desempeño de forma objetiva. Por ello,
en el siguiente apartado se describen los \textit{benchmarks} y
conjuntos de datos (\textit{datasets}) estándar en la industria que se
emplean para la validación experimental de los métodos propuestos en el
presente trabajo.

\section{Datasets y benchmarks}
\label{sec:datasets}

La validación objetiva de algoritmos de percepción para conducción
autónoma requiere el uso de entornos de evaluación estandarizados que
proporcionen \textit{ground truth} preciso y permitan comparar de forma
cuantitativa el rendimiento de distintos métodos en condiciones
reproducibles.

Las métricas más utilizadas en estos \textit{benchmarks} incluyen la
precisión media (\textit{mean Average Precision}, mAP) para tareas de
detección de objetos, la Intersección sobre Unión (\textit{Intersection
over Union}, IoU) para segmentación semántica y evaluación de cajas
delimitadoras, el \textbf{F1-\textit{score}} como medida armónica entre
precisión y \textit{recall}, y la \textbf{latencia por frame} como
indicador de viabilidad en tiempo real.

A continuación, se describen los \textit{benchmarks} más relevantes
para tareas de segmentación y detección de obstáculos basadas en LiDAR.

\subsection{SemanticKITTI}
\label{sec:semantickitti}

SemanticKITTI~\cite{behley2019semantickitti} es un \textit{dataset} de
gran escala construido sobre las secuencias de odometría del
\textit{dataset} original KITTI~\cite{geiger2012kitti}, al que añade
anotaciones semánticas densas punto a punto para cada nube LiDAR. Las
nubes fueron capturadas con un Velodyne HDL-64E (64 canales, 10~Hz)
montado sobre un vehículo Volkswagen Passat circulando por la ciudad de
Karlsruhe (Alemania), en condiciones urbanas convencionales y diurnas.

El \textit{dataset} contiene 22 secuencias numeradas del 00 al 21, con
un total de 43.552 \textit{frames} anotados según 28 clases semánticas
que comprenden desde elementos estructurales (calzada, acera,
edificios, vegetación) hasta objetos dinámicos (coches, peatones,
ciclistas, camiones) y elementos auxiliares (postes, señales de
tráfico, vallas). La convención estándar establecida por los autores
reserva las secuencias 00--07 y 09--10 para entrenamiento, la secuencia
08 para validación y las secuencias 11--21 para evaluación oficial en
el \textit{benchmark online}, cuyas etiquetas de \textit{ground truth}
no son públicas para evitar \textit{overfitting} al conjunto de prueba.

SemanticKITTI se ha consolidado como el estándar para la evaluación de
métodos de segmentación semántica 3D sobre LiDAR en entornos urbanos, y
es uno de los \textit{benchmarks} más citados en la literatura reciente
de percepción para vehículos autónomos.

\subsection{3D-Curb}
\label{sec:3dcurb}

3D-Curb~\cite{zhao2024curbnet} es un conjunto de anotaciones
especializado en la detección de bordillos construido como extensión de
SemanticKITTI. Partiendo de las mismas secuencias LiDAR del
\textit{dataset} original, añade una clase específica: \textit{curb}.

La aportación fundamental de 3D-Curb es proporcionar un
\textit{benchmark point-wise} para una tarea que SemanticKITTI no
aborda explícitamente: la discriminación geométrica entre calzada,
acera y el propio bordillo que los separa. Dado que los bordillos son
estructuras de baja altura y quedan habitualmente absorbidos dentro de
la clase \textit{road} o \textit{sidewalk} en otras anotaciones,
3D-Curb permite evaluar cuantitativamente tanto métodos geométricos
clásicos como arquitecturas de aprendizaje profundo específicamente
diseñadas para detección de bordillos. Las métricas habituales
reportadas sobre este \textit{dataset} son la precisión, el
\textit{recall} y el F1-\textit{score} específicos de la clase
\textit{curb}, calculados únicamente sobre los puntos anotados como
bordillo.

\subsection{KITTI-360}
\label{sec:kitti360}

KITTI-360~\cite{liao2023kitti360} representa la evolución del
\textit{dataset} KITTI clásico, ampliando significativamente la
cobertura geográfica y la riqueza de las anotaciones. Ofrece una visión
de 360 grados del entorno mediante la integración de cámaras
omnidireccionales y LiDAR de alta resolución.

La contribución distintiva de KITTI-360 reside en su esquema de
anotación denso y temporalmente consistente: además de etiquetas
semánticas 2D y 3D, el \textit{dataset} proporciona \textit{bounding
boxes} 3D anotadas manualmente (no derivadas de segmentación panóptica
como en otros \textit{datasets}) para objetos estáticos y dinámicos a
lo largo de secuencias temporales largas, junto con máscaras de
instancia que mantienen la consistencia del identificador de objeto
entre \textit{frames} consecutivos. Esta característica lo convierte en
uno de los pocos \textit{benchmarks} públicos que permite evaluar
rigurosamente la calidad geométrica de las OBBs predichas por un
sistema de percepción, sin depender de anotaciones automáticas
derivadas de segmentación semántica.

\subsection{German Outdoor and Offroad Dataset (GOOSE)}
\label{sec:goose}

GOOSE~\cite{mortimer2024goose} surge como respuesta a la falta de
\textit{benchmarks} específicos para entornos \textit{off-road} y no
estructurados, un dominio todavía poco representado en la literatura de
percepción basada en LiDAR, tradicionalmente centrada en escenarios
urbanos. El conjunto de datos fue adquirido utilizando una plataforma
de investigación militar, el vehículo MuCAR-3 (un Volkswagen Touareg
modificado por la Universidad de las Fuerzas Armadas de Múnich),
equipada con múltiples sensores: un Velodyne Alpha Prime (VLS-128, 128
canales, 10~Hz, FOV vertical de 40\textdegree) montado sobre el techo,
dos LiDAR Ouster OS0 laterales inclinados $\pm 45^{\circ}$ para
detección próxima al chasis, cámaras RGB+NIR, cámara térmica infrarroja
y sistema de posicionamiento inercial.

El \textit{dataset} proporciona 10.000 \textit{frames} LiDAR y 10.000
\textit{frames} de imagen anotados con 64 clases semánticas
específicamente diseñadas para entornos \textit{off-road}, agrupadas en
categorías como vegetación alta y baja, distintos tipos de terreno
(asfalto, grava, tierra, nieve, hierba), estructuras artificiales,
vehículos militares y civiles, y elementos de señalización rural. La
diversidad de escenarios cubre ocho entornos distintos (bosques densos,
caminos agrícolas, zonas suburbanas, terrenos con vegetación variable)
recogidos a lo largo de las cuatro estaciones del año y bajo
condiciones meteorológicas cambiantes.
